\begin{proof}
    First, it must be shown that $\Sigma_{\CatHQC} = \bigsqcup_{n=0}^\infty \Sigma_n^*$.
    This will require showing that if $j \ne k$, then $\Sigma_j^*$ is disjoint from $\Sigma_k^*$.
    To this end, let $j \in \mathbb{N}$ and $k \in \mathbb{N}$ such that $j \ne k$.
    Assume, without loss of generality, that $j < k$.
    There are two cases to consider.
    \begin{enumerate}
    \item If $j = 0$, then $\Sigma_j^* = \Sigma_{\CatHQC}^j \subseteq \Sigma_{\CatHQC}^j$.
    \item If $j > 0$, then $\Sigma_j^* = \Sigma_{\CatHQC}^j \setminus \Sigma_{\CatHQC}^{j-1} \subseteq \Sigma_{\CatHQC}^j$.
    \end{enumerate}
    In either case, $\Sigma_j^* \subseteq \Sigma_{\CatHQC}^j$.
    Since $j \le k - 1$, then by the inductive construction $\Sigma_j^* \subseteq \Sigma_{\CatHQC}^{k-1}$.
    Since $\Sigma_k^* = \Sigma_{\CatHQC}^k \setminus \Sigma_{\CatHQC}^{k-1}$, then $\Sigma_j^* \cap \Sigma_k^* = \varnothing$.
    Since $j$ and $k$ were arbitrary, then $\Sigma_j^*$ and $\Sigma_k^*$ are disjoint for each $j \in \mathbb{N}$ and $k \in \mathbb{N}$.
    Then $\bigsqcup_{n=0}^\infty \Sigma_n^* \subseteq \bigcup_{n=0}^\infty \Sigma_{\CatHQC}^n = \Sigma_{\CatHQC}$.
    Then let $g \in \Sigma_{\CatHQC}$.
    Then there exists some least $n \in \mathbb{N}$ such that $g \in \Sigma_{\CatHQC}^n$.
    There are two cases to consider.
    \begin{enumerate}
    \item If $n = 0$, then $g \in \Sigma_{\CatHQC}^n = \Sigma_n^*$.
    \item If $n > 0$, then $g \not \in \Sigma_{\CatHQC}^{n-1}$.
          Then $g \in \Sigma_{\CatHQC}^n \setminus \Sigma_{\CatHQC}^{n-1} = \Sigma_n^*$.
    \end{enumerate}
    In either case, $g \in \Sigma_n^*$.
    Since $g$ was arbitrary, then $\Sigma_{\CatHQC} = \bigsqcup_{n=0}^\infty \Sigma_n^*$.
    Then for each $g \in \Sigma_{\CatHQC}$, there exists a unique $\nu(g) \in \mathbb{N}$ such that $g \in \Sigma_{\nu(g)}^*$.
    This defines an assignment $\nu: \Sigma_{\CatHQC} \to \mathbb{N}$.
    This define a prop signature morphism $\rho: \Sigma_{\CatHQC} \to U( \mathcal{C} )$ such that $\rho( g ) = \tau_{\nu(g)}( g )$.
    This is well-defined since $g \in \Sigma_{\nu(g)}^*$ and $\tau_{\nu(g)}: \Sigma_{\nu(g)}^* \to U( \mathcal{C} )$.
    Then there exists a unique prop functor $F: \CatHQC \to \mathcal{C}$ such that $F( g ) = \rho( g )$ for each $g \in \Sigma_{\CatHQC}$.
    Moreover, if $n \in \mathbb{N}$ and $g \in \Sigma_n^*$, then $\nu( g ) = n$ and consequently $F( g ) = \rho( g ) = \tau_{\nu(g)}( g ) = \tau_n( g )$.
    It remains to be shown that $F$ is unique.
    Assume that $G: \CatHQC \to \mathcal{C}$ is a prop functor such that $G( g ) = \tau_n( g )$ for each $n \in \mathbb{N}$ and $g \in \Sigma_n^*$.
    Let $g \in \Sigma_{\CatHQC}$.
    Then $g \in \Sigma_{\nu(g)}^*$.
    Then $G( g ) = \tau_{\nu(g)}( g ) = F( g )$.
    Since $g$ was arbitrary, then $G \circ i = F \circ i$.
    Then by the universal property of free prop categories, $G = F$.
    Since $G$ was arbitrary, then $F$ is unique.
\end{proof}